% IEEE Journal Template
% Usage: Replace placeholders with your content
% Compile with: pdflatex -> bibtex -> pdflatex -> pdflatex

\documentclass[journal]{IEEEtran}

\usepackage{amsmath,amssymb,amsfonts}
\usepackage{algorithmic}
\usepackage{graphicx}
\usepackage{textcomp}
\usepackage{xcolor}
\usepackage{cite}
\usepackage{hyperref}
\usepackage{booktabs}
\usepackage{multirow}

\begin{document}

\title{TODO: Paper Title}

\author{
  \IEEEauthorblockN{TODO: First Author\IEEEauthorrefmark{1},
  TODO: Second Author\IEEEauthorrefmark{2}}
  \IEEEauthorblockA{
    \IEEEauthorrefmark{1}TODO: Department, University, Country\\
    Email: todo@university.edu
  }
  \IEEEauthorblockA{
    \IEEEauthorrefmark{2}TODO: Department, University, Country\\
    Email: todo@university.edu
  }
}

\maketitle

\begin{abstract}
TODO: Write abstract (150-250 words).
Structure: Problem -- Gap -- Method -- Key Result -- Implication.
\end{abstract}

\begin{IEEEkeywords}
TODO: keyword1, keyword2, keyword3, keyword4, keyword5
\end{IEEEkeywords}

% ============================================================
\section{Introduction}
\label{sec:introduction}

TODO: Motivation, gap, approach, main findings, contributions.

\subsection{Contributions}
The main contributions of this paper are:
\begin{itemize}
  \item TODO: Contribution 1 (specific, evidence-backed)
  \item TODO: Contribution 2
  \item TODO: Contribution 3
\end{itemize}

% ============================================================
\section{Related Work}
\label{sec:related_work}

\subsection{TODO: Cluster 1}
TODO: Group related papers by theme, end with contrast to your work.

\subsection{TODO: Cluster 2}
TODO: Second thematic group.

% ============================================================
\section{Method}
\label{sec:method}

\subsection{Problem Formulation}
TODO: Problem setup, notation, assumptions.

\subsection{Proposed Approach}
TODO: Model/algorithm description, design choices, training procedure.

% ============================================================
\section{Experiments}
\label{sec:experiments}

\subsection{Experimental Setup}
TODO: Datasets, baselines, metrics, implementation details.

\subsection{Main Results}
TODO: Results table, comparison with baselines.

\begin{table}[htbp]
\caption{TODO: Main Results}
\label{tab:main_results}
\centering
\begin{tabular}{lccc}
\toprule
Method & Metric 1 & Metric 2 & Params \\
\midrule
Baseline 1 & TODO & TODO & TODO \\
Baseline 2 & TODO & TODO & TODO \\
\textbf{Ours} & \textbf{TODO} & \textbf{TODO} & TODO \\
\bottomrule
\end{tabular}
\end{table}

\subsection{Ablation Studies}
TODO: Ablation results showing each component's contribution.

% ============================================================
\section{Discussion}
\label{sec:discussion}
TODO: Interpret results, discuss implications, compare with expectations.

% ============================================================
\section{Limitations}
\label{sec:limitations}
TODO: Scope limits, dataset limits, compute limits, generalization risks.

% ============================================================
\section{Conclusion}
\label{sec:conclusion}
TODO: Restate contribution and strongest evidence. No new claims.

% ============================================================
\section*{Acknowledgments}
TODO: Funding, support, compute resources.

\bibliographystyle{IEEEtran}
\bibliography{references}

\end{document}
